% TeX root=../main.tex

\begin{table}[!htb]
	\caption{Sistema consonantal do indo-europeu ao \glsxtrlong{eslavoantigo}}\label{tab:pieeslavcons}
	\begin{center}
		\begin{tabular}[c]{lll}
			\toprule
			*\emph{p} > \emph{p}                      &                             & *\pietrans{bh} > \emph{b}    \\
			\rowcolor{gray!10} *\emph{t} > \emph{t}   & *\emph{d} > \emph{d}        & *\pietrans{dh} > \ocsex{d}   \\
			*\pietrans{c} > \emph{s}                  & *\pietrans{j} > \ocsex{z}   & *\pietrans{jh} > \ocsex{z}   \\
			\rowcolor{gray!10} *\emph{k} > \emph{k/č} & *\emph{g} > \emph{g/ž}      & *\pietrans{gh} > \ocsex{g/ž} \\
			*\pietrans{kw} > \emph{k/č}               & *\pietrans{gw} > \emph{g/ž} & *\pietrans{gwh} > \emph{g/ž} \\
			\rowcolor{gray!10} \multicolumn{3}{l}{*\emph{r}, *\emph{l}, *\emph{n}, *\emph{m} > \emph{r},
			\emph{l}, \emph{n}, \emph{m}}                                                                          \\
			\multicolumn{3}{l}{*\emph{s} > \emph{s/š/x} }                                                          \\
			\multicolumn{3}{l}{de acordo com o contexto (RUKI)}                                                    \\
			\rowcolor{gray!10} \multicolumn{3}{l}{*\pietrans{y}̯, *\pietrans{w} > \ocsex{j}, \ocsex{v}}             \\
			\bottomrule
			\multicolumn{3}{l}{}                                                                                   \\
			\multicolumn{3}{l}{NB\@: o indo-europeu não possuía /b/.}                                              \\
		\end{tabular}
	\end{center}
\end{table}
